\begin{textsample}{POS Dim 3 – human – Score 5.00 – rj/rj\_ef\_intro\_4\_p5.txt}  \label{ex:f3_pos_013}
O \textbf{Documento} de Orientação \textbf{Curricular} aqui apresentado busca tratar, o estudante, como cidadão de direito, que independente de história, origem, raça, credo, cor, tem garantidos legalmente direitos para todos.
Que \textbf{proposta} \textbf{curricular} se faz para o estudante do Rio de Janeiro?
Que \textbf{proposta} \textbf{curricular} se pensa para as múltiplas realidades espalhados nos 92 municípios?
Que tessitura se faz com docentes que se debruçam na criação do cotidiano?
Como dialogar com componentes \textbf{curriculares} sem perder a originalidade e autoria regional do que se pensa acerca do sujeito e educação?

% matched lemmas: curricular, documento, proposta
\end{textsample}
